% ============================================================
% ch2.tex — 2 МАТЕМАТИЧЕСКАЯ МОДЕЛЬ (этап 1)
% ============================================================
\chapter{Математическая модель}
\label{ch:model}

Математическая модель задачи состоит из дифференциального уравнения,
которому удовлетворяет температура внутри цилиндра, и условий на его
границе. Ниже перечислены принятые допущения, выведены уравнение и
граничные условия, записана математическая постановка задачи и рассмотрены
её свойства. При построении модели использованы материалы лекций курса
и учебная литература~[1--3].

\section{Принятые допущения}
\label{sec:assumptions}

Модель строится при следующих допущениях:
\begin{enumerate}
  \item цилиндр сплошной, его материал однороден и изотропен: коэффициент
        теплопроводности одинаков во всех точках и по всем направлениям,
        а также не зависит от температуры;
  \item внутри цилиндра нет источников и стоков тепла;
  \item тепло внутри цилиндра переносится только теплопроводностью;
  \item тепловой режим стационарный: температура в каждой точке не зависит
        от времени;
  \item торцы идеально теплоизолированы, а температура боковой поверхности
        в точности равна заданной и на высоте $z=c$ меняется скачком.
\end{enumerate}

\section{Уравнение стационарной теплопроводности}
\label{sec:heat}

Пусть $T(M,t)$~— температура в точке~$M$ тела в момент времени~$t$.
Перенос тепла теплопроводностью описывается законом Фурье~[2,\,3]:
\begin{equation}
\label{eq:fourier}
\vec q=-k\,\Grad T,
\postdisplaypenalty=10000
\end{equation}
где $\vec q$~— вектор плотности теплового потока, Вт/м$^2$;\\
\hphantom{где }$k$~— коэффициент теплопроводности, Вт/(м$\cdot$К).

Вектор~$\vec q$ направлен в сторону переноса тепла, а его модуль равен
количеству тепла, проходящему за единицу времени через единичную площадку,
перпендикулярную потоку. Знак «минус» показывает, что тепло переносится
в сторону убывания температуры.

Выделим внутри тела произвольный объём~$V$, ограниченный замкнутой
поверхностью~$S$ с внешней нормалью~$\vec n$, и запишем для него баланс
тепла за единицу времени:
\begin{equation}
\label{eq:balance}
\iiint\limits_V C\rho\,\pd{T}{t}\,dV
=-\iint\limits_S \vec q\cdot\vec n\,dS+\iiint\limits_V F\,dV,
\postdisplaypenalty=10000
\end{equation}
где $C$~— удельная теплоёмкость, Дж/(кг$\cdot$К);\\
\hphantom{где }$\rho$~— плотность, кг/м$^3$;\\
\hphantom{где }$F$~— мощность внутренних источников тепла в единице объёма,
Вт/м$^3$.

Левая часть равенства~\eqref{eq:balance}~— скорость изменения количества
тепла, запасённого в объёме~$V$. Первое слагаемое правой части~— тепло,
поступающее в объём через поверхность~$S$, второе~— тепло, выделяемое
источниками.

По формуле Остроградского поток вектора~$\vec q$ через замкнутую поверхность
равен интегралу от его расходимости по объёму:
\begin{equation}
\label{eq:ostrogradsky}
\iint\limits_S \vec q\cdot\vec n\,dS=\iiint\limits_V \Div\vec q\,dV .
\end{equation}
Заменяя по этой формуле поверхностный интеграл в~\eqref{eq:balance}
объёмным и подставляя закон Фурье~\eqref{eq:fourier}, получаем
\begin{equation}
\label{eq:balance-vol}
\iiint\limits_V\Bigl(C\rho\,\pd{T}{t}-\Div\bigl(k\,\Grad T\bigr)-F\Bigr)dV=0 .
\end{equation}
Объём~$V$ произволен, а подынтегральная функция непрерывна, поэтому она
равна нулю в каждой точке тела:
\begin{equation}
\label{eq:heat-general}
C\rho\,\pd{T}{t}=\Div\bigl(k\,\Grad T\bigr)+F .
\end{equation}

Для однородного тела $k=\mathrm{const}$ и
$\Div\bigl(k\,\Grad T\bigr)=k\,\Delta T$, где $\Delta$~— оператор Лапласа:
\begin{equation}
\label{eq:lap-cart}
\Delta T=\pdd{T}{x}+\pdd{T}{y}+\pdd{T}{z} .
\end{equation}
Из~\eqref{eq:heat-general} получаем уравнение теплопроводности
\begin{equation}
\label{eq:heat}
\pd{T}{t}=\chi\,\Delta T+\frac{F}{C\rho},
\qquad \chi=\frac{k}{C\rho},
\postdisplaypenalty=10000
\end{equation}
где $\chi$~— коэффициент температуропроводности, м$^2$/с. В лекциях курса
это уравнение записано в виде $u_t=a^2\Delta u+f$, то есть $\chi$
соответствует $a^2$, а $F/(C\rho)$~— функции~$f$. Буквы $a$ и $c$ в условии
задачи уже использованы (радиус цилиндра и высота, на которой температура
боковой поверхности меняется скачком), поэтому коэффициент
температуропроводности обозначен~$\chi$, а удельная теплоёмкость~— $C$.

Режим стационарный, поэтому $\partial T/\partial t=0$; источников тепла
нет, поэтому $F=0$. Уравнение~\eqref{eq:heat} переходит в уравнение Лапласа
\begin{equation}
\label{eq:laplace}
\Delta T=0 .
\end{equation}
Уравнение~\eqref{eq:laplace} не содержит характеристик материала $k$, $C$
и~$\rho$. Следовательно, стационарное распределение температуры
в однородном теле без внутренних источников определяется только формой тела
и условиями на его границе.

\section{Цилиндрические координаты и осевая симметрия}
\label{sec:cyl}

Границу цилиндра удобно описывать в цилиндрических координатах
$(r,\varphi,z)$, которые связаны с декартовыми координатами формулами
\begin{equation}
\label{eq:cyl}
x=r\cos\varphi,\qquad y=r\sin\varphi,\qquad z=z,
\postdisplaypenalty=10000
\end{equation}
где $r\ge 0$~— расстояние от точки до оси~$Oz$;\\
\hphantom{где }$\varphi$~— полярный угол, $0\le\varphi<2\pi$.

Ось~$Oz$ совпадает с осью цилиндра и направлена от нижнего торца
к верхнему, начало координат помещено в центр нижнего торца. В этих
координатах цилиндр занимает область
\begin{equation}
\label{eq:Omega}
\Omega=\bigl\{(r,\varphi,z):\ 0\le r<a,\ \ 0\le\varphi<2\pi,\ \ 0<z<h\bigr\},
\end{equation}
боковая поверхность задаётся уравнением $r=a$, а торцы~— уравнениями
$z=0$ и $z=h$.

Оператор Лапласа в цилиндрических координатах имеет вид~[3]
\begin{equation}
\label{eq:lap-cyl}
\Delta T=\frac{1}{r}\,\pd{}{r}\Bigl(r\,\pd{T}{r}\Bigr)
+\frac{1}{r^2}\,\pdd{T}{\varphi}+\pdd{T}{z} .
\end{equation}

Цилиндр~— тело вращения с осью~$Oz$, а условия на его поверхности не
зависят от угла~$\varphi$: заданная условием~\eqref{eq:task} температура
боковой поверхности зависит только от~$z$, торцы теплоизолированы целиком.
Ни одно значение угла~$\varphi$ не выделено, поэтому искомое температурное
поле осесимметрично:
\begin{equation}
\label{eq:axisym}
T=T(r,z),\qquad \pd{T}{\varphi}=0 .
\end{equation}
Строгое обоснование равенств~\eqref{eq:axisym} опирается на единственность
решения и приведено в подразделе~\ref{sec:type}.
Подставляя~\eqref{eq:lap-cyl} и~\eqref{eq:axisym} в уравнение
Лапласа~\eqref{eq:laplace}, получаем
\begin{equation}
\label{eq:main-div}
\frac{1}{r}\,\pd{}{r}\Bigl(r\,\pd{T}{r}\Bigr)+\pdd{T}{z}=0,
\end{equation}
или, после дифференцирования произведения,
\begin{equation}
\label{eq:main}
\pdd{T}{r}+\frac{1}{r}\,\pd{T}{r}+\pdd{T}{z}=0 .
\end{equation}
Задача стала двумерной. Функция $T(r,z)$ ищется в прямоугольнике
\begin{equation}
\label{eq:Pi}
\Pi=\bigl\{(r,z):\ 0<r<a,\ \ 0<z<h\bigr\},
\end{equation}
который представляет собой половину осевого сечения цилиндра.

\section{Граничные условия}
\label{sec:bc}

\emph{Боковая поверхность.} По условию боковая поверхность поддерживается
при заданной температуре, то есть в каждой её точке известно значение
искомой функции:
\begin{equation}
\label{eq:bc-side}
T(a,z)=\mu(z),\quad 0<z<h,
\qquad
\mu(z)=
\begin{cases}
0,   & 0<z<c,\\[2pt]
T_0, & c\le z<h .
\end{cases}
\end{equation}
Это граничное условие I~рода (условие Дирихле). Оно неоднородно, а функция
$\mu(z)$ кусочно-постоянна и имеет разрыв первого рода в точке $z=c$.

\emph{Торцы.} Количество тепла, проходящее за единицу времени через
единицу площади поверхности тела в направлении внешней нормали~$\vec n$,
в силу закона Фурье~\eqref{eq:fourier} равно
\begin{equation}
\label{eq:flux-n}
q_n=\vec q\cdot\vec n=-k\,\Grad T\cdot\vec n=-k\,\pd{T}{n} .
\end{equation}
Торец теплоизолирован, если тепло через него не проходит: $q_n=0$, то есть
$\partial T/\partial n=0$. На верхнем торце $z=h$ внешняя нормаль направлена
вдоль оси~$Oz$ и $\partial T/\partial n=\partial T/\partial z$. На нижнем
торце $z=0$ она направлена противоположно оси и
$\partial T/\partial n=-\partial T/\partial z$. В обоих случаях получаем
\begin{equation}
\label{eq:bc-ends}
\pd{T}{z}\bigg|_{z=0}=0,\qquad \pd{T}{z}\bigg|_{z=h}=0,\qquad 0\le r<a .
\end{equation}
Это однородные граничные условия II~рода (условия Неймана).

\emph{Ось цилиндра.} Отрезок $r=0$ не является частью поверхности
цилиндра: точки оси лежат внутри тела. Однако он служит стороной
прямоугольника~$\Pi$, а коэффициент $1/r$ уравнения~\eqref{eq:main} при
$r\to 0$ неограниченно растёт. У~этого уравнения есть решения, которые на
оси обращаются в бесконечность, например $T=A\ln(r/a)$, где $A\ne 0$~—
постоянная. Такое решение соответствует источнику или стоку тепла на оси,
а по принятым допущениям их нет. Поэтому на оси ставится условие
ограниченности
\begin{equation}
\label{eq:bc-axis}
\bigl|T(0,z)\bigr|<\infty,\qquad 0\le z\le h .
\end{equation}
Для осесимметричного поля, гладкого на оси, из симметрии следует также
равенство
\begin{equation}
\label{eq:bc-axis-n}
\pd{T}{r}\bigg|_{r=0}=0,
\end{equation}
которое удобно использовать при численном решении.

\section{Математическая постановка задачи}
\label{sec:statement}

До использования осевой симметрии задача ставится в трёхмерной
области~$\Omega$: требуется найти функцию $T(r,\varphi,z)$ из условий
\begin{equation}
\label{eq:problem3d}
\begin{cases}
\Delta T=0,
  & 0<r<a,\ \ 0\le\varphi<2\pi,\ \ 0<z<h,\\[6pt]
\pd{T}{z}\bigg|_{z=0}=0,\qquad \pd{T}{z}\bigg|_{z=h}=0,
  & 0\le r<a,\ \ 0\le\varphi<2\pi,\\[10pt]
T\big|_{r=a}=\mu(z),
  & 0\le\varphi<2\pi,\ \ 0<z<h,\\[6pt]
\bigl|T\big|_{r=0}\bigr|<\infty,
  & 0\le z\le h,\\[6pt]
T(r,\varphi+2\pi,z)=T(r,\varphi,z),
  & 0<r\le a,\ \ 0\le z\le h,
\end{cases}
\postdisplaypenalty=10000
\end{equation}
где оператор Лапласа имеет вид~\eqref{eq:lap-cyl}, а последнее условие~—
условие периодичности по углу: значениям $\varphi$ и $\varphi+2\pi$
отвечает одна и та же точка. Граничные условия
задачи~\eqref{eq:problem3d} не зависят от~$\varphi$, поэтому её решение
осесимметрично: слагаемое с $\partial^2T/\partial\varphi^2$ исчезает,
условие периодичности выполняется автоматически, и задача сводится
к двумерной.

Объединяя уравнение~\eqref{eq:main} с условиями
\eqref{eq:bc-side}, \eqref{eq:bc-ends}, \eqref{eq:bc-axis}, приходим
к итоговой постановке. Требуется найти функцию $T(r,z)$, дважды непрерывно
дифференцируемую в прямоугольнике~$\Pi$ и удовлетворяющую условиям
\begin{equation}
\label{eq:problem}
\begin{cases}
\pdd{T}{r}+\dfrac{1}{r}\,\pd{T}{r}+\pdd{T}{z}=0,
  & 0<r<a,\ \ 0<z<h,\\[10pt]
\pd{T}{z}\bigg|_{z=0}=0,\qquad \pd{T}{z}\bigg|_{z=h}=0,
  & 0\le r<a,\\[10pt]
T\big|_{r=a}=\mu(z),
  & 0<z<h,\\[6pt]
\bigl|T\big|_{r=0}\bigr|<\infty,
  & 0\le z\le h,
\end{cases}
\end{equation}
где $T(r,z)$~— искомая функция;\\
\hphantom{где }$a$, $h$, $c$, $T_0$~— заданные параметры;\\
\hphantom{где }$\mu(z)$~— известная функция: $\mu(z)=0$ при
\mbox{$0<z<c$} и $\mu(z)=T_0$ при \mbox{$c\le z<h$}
(формула~\eqref{eq:bc-side}).

В лекциях курса решение краевой задачи ищется в классе функций, дважды
непрерывно дифференцируемых в области и непрерывных в замкнутой области.
В задаче~\eqref{eq:problem} граничная функция~$\mu(z)$ разрывна, и решение
не может быть непрерывным в точке $(a,c)$. Поэтому вместо непрерывности
в замкнутой области от него требуются ограниченность в области~$\Pi$
и непрерывность вплоть до границы всюду, кроме этой точки. Значение
функции~$\mu$ в самой точке $z=c$ на решение не влияет. Чтобы условия на
торцах имели смысл, производная $\partial T/\partial z$ должна быть
непрерывна вплоть до отрезков $z=0$ и $z=h$, $0\le r<a$.

В лекциях граница области предполагается гладкой, а граница цилиндра лишь
кусочно-гладкая: на рёбрах $r=a$, $z=0$ и $r=a$, $z=h$ нормаль не
определена. Поэтому условия~\eqref{eq:bc-ends} поставлены при $r<a$,
условие~\eqref{eq:bc-side}~— при $0<z<h$, а на самих рёбрах условия не
задаются. Расчётная область и условия на её границе показаны на
рисунке~\ref{fig:domain}.

\begin{figure}[H]
  \centering
  % Рисунок 2 — расчётная область (меридиональное сечение цилиндра)
\begin{tikzpicture}[>={Stealth[length=2.6mm]}, line cap=round, line join=round]
  \def\A{5.6}
  \def\H{6.2}
  \def\C{2.5}

  % --- теплоизоляция торцов (штриховка снаружи области) ---
  \fill[pattern=north east lines, pattern color=black!55] (0,-0.28) rectangle (\A,0);
  \fill[pattern=north east lines, pattern color=black!55] (0,\H) rectangle (\A,\H+0.28);

  % --- область ---
  \fill[gray!10] (0,0) rectangle (\A,\H);
  \draw[very thick] (0,0) -- (\A,0);
  \draw[very thick] (0,\H) -- (\A,\H);
  \draw[very thick] (\A,0) -- (\A,\C);
  \draw[line width=2.2pt, gray!70!black] (\A,\C) -- (\A,\H);
  \draw[dash dot, thick] (0,0) -- (0,\H);

  % --- оси ---
  \draw[->] (0,\H) -- (0,\H+1.35) node[left] {$z$};
  \draw[->] (\A,0) -- (\A+1.9,0) node[below] {$r$};
  \draw[dash dot] (0,0) -- (0,-0.9);

  % --- отметки ---
  \node[below left=-1pt] at (0,0) {$O$};
  \node[below=6pt] at (\A,0) {$a$};
  \draw[thin] (-0.12,\C) -- (0.12,\C);
  \node[left=3pt] at (0,\C) {$c$};
  \node[left=3pt] at (0,\H) {$h$};
  \draw[thin, dashed] (0,\C) -- (\A,\C);
  \fill (\A,\C) circle[radius=2pt];

  % --- уравнение внутри области ---
  \node at (0.5*\A, 0.5*\C+0.5*\H+0.25)
    {$\dfrac{1}{r}\,\pd{}{r}\!\Bigl(r\,\pd{T}{r}\Bigr)+\pdd{T}{z}=0$};
  \node at (0.5*\A, 0.5*\C) {$\Pi$};

  % --- граничные условия ---
  \node[below] at (0.5*\A,-0.32) {$\pd{T}{z}=0$};
  \node[above] at (0.5*\A,\H+0.3) {$\pd{T}{z}=0$};
  \node[right=4pt] at (\A, 0.5*\C+0.2) {$T=0$};
  \node[right=4pt] at (\A, 0.5*\C+0.5*\H) {$T=T_0$};
  \node[left=4pt] at (0, 0.5*\C+0.5*\H+0.25) {$|T|<\infty$};

  % --- внешние нормали к поверхности цилиндра ---
  \draw[->, thick] (0.78*\A,\H) -- (0.78*\A,\H+0.95) node[right] {$\vec n$};
  \draw[->, thick] (0.78*\A,0) -- (0.78*\A,-0.95) node[right] {$\vec n$};
  \draw[->, thick] (\A,0.9*\H) -- (\A+0.95,0.9*\H) node[above] {$\vec n$};
\end{tikzpicture}

  \caption{Расчётная область и граничные условия}
  \label{fig:domain}
\end{figure}

\section{Тип уравнения и свойства задачи}
\label{sec:type}

\emph{Тип уравнения.} Уравнение~\eqref{eq:main}~— линейное однородное
уравнение второго порядка с двумя независимыми переменными $r$ и~$z$.
Сравним его с общим видом уравнения, линейного относительно старших
производных~[2], записанным для независимых переменных $r$, $z$ и искомой
функции~$T$:
\begin{equation}
\label{eq:general}
a_{11}\,\pdd{T}{r}+2a_{12}\,\frac{\partial^2 T}{\partial r\,\partial z}
+a_{22}\,\pdd{T}{z}+\Phi\Bigl(r,z,T,\pd{T}{r},\pd{T}{z}\Bigr)=0 .
\end{equation}
Для уравнения~\eqref{eq:main} $a_{11}=1$, $a_{12}=0$, $a_{22}=1$,
$\Phi=(1/r)\,\partial T/\partial r$ (коэффициенты $a_{11}$, $a_{12}$,
$a_{22}$ не связаны с радиусом цилиндра~$a$). Тогда
\begin{equation}
\label{eq:discr}
D=a_{12}^2-a_{11}a_{22}=-1<0,
\end{equation}
и уравнение~\eqref{eq:main} относится к эллиптическому типу во всей
области~$\Pi$. Слагаемое с первой производной на тип уравнения не влияет.

Величина~$D$~— дискриминант характеристического уравнения
$a_{11}\,dz^2-2a_{12}\,dr\,dz+a_{22}\,dr^2=0$. Для
уравнения~\eqref{eq:main} оно имеет вид $dz^2+dr^2=0$ и вещественных решений
не имеет. Значит, вещественных характеристик у уравнения~\eqref{eq:main}
нет, а его старшая часть~— сумма вторых производных по $r$ и по~$z$~— уже
записана в каноническом виде. Исходное уравнение Лапласа~\eqref{eq:laplace}
с тремя независимыми переменными также относится к эллиптическому типу: все
коэффициенты его характеристической квадратичной формы
$\lambda_1^2+\lambda_2^2+\lambda_3^2$ отличны от нуля и имеют один знак.

\emph{Тип задачи.} Время в уравнение не входит, начальных условий нет,
поэтому~\eqref{eq:problem}~— краевая задача. На разных частях границы
заданы условия разного рода: на боковой поверхности~— I~рода, на торцах~—
II~рода. При этом задача не является смешанной в том смысле, который принят
в курсе лекций: смешанной там называется начально-краевая задача. Уравнение
и условия на торцах однородны, единственная неоднородность задачи~—
функция~$\mu(z)$ в условии на боковой поверхности.

\emph{Единственность решения.} Пусть $T_1$ и $T_2$~— два решения
трёхмерной задачи~\eqref{eq:problem3d}. Их разность $w=T_1-T_2$
удовлетворяет уравнению Лапласа $\Delta w=0$ в области~$\Omega$, условию
${w=0}$ на боковой поверхности и условию ${\partial w/\partial n=0}$ на
торцах. Так как
$\Div(w\,\Grad w)=|\Grad w|^2+w\,\Delta w=|\Grad w|^2$, формула
Остроградского~\eqref{eq:ostrogradsky} для вектора $w\,\Grad w$ и
области~$\Omega$ даёт
\begin{equation}
\label{eq:energy}
\iiint\limits_{\Omega}|\Grad w|^{2}\,dV
=\iint\limits_{S_\Omega} w\,\pd{w}{n}\,dS=0,
\end{equation}
где $S_\Omega$~— поверхность цилиндра. Поверхностный интеграл равен нулю,
потому что на боковой поверхности ${w=0}$, а на торцах
${\partial w/\partial n=0}$. Подынтегральная функция слева неотрицательна,
поэтому $\Grad w=0$ всюду в~$\Omega$, то есть $w=\mathrm{const}$, а из
условия на боковой поверхности следует ${w\equiv 0}$. Значит,
$T_1\equiv T_2$: задача~\eqref{eq:problem3d} не может иметь двух различных
решений. Тем самым единственно и решение двумерной
задачи~\eqref{eq:problem}: каждое её решение~— это не зависящее
от~$\varphi$ решение задачи~\eqref{eq:problem3d}.

Равенство~\eqref{eq:energy} справедливо, если первые производные
функции~$w$ ограничены вплоть до поверхности цилиндра. В рассматриваемой
задаче функция~$\mu(z)$ разрывна, и о решениях $T_1$ и $T_2$ вблизи
окружности $r=a$, $z=c$ известно лишь то, что они ограничены;
ограниченность производных их разности заранее не гарантирована. В этом
случае единственность ограниченного решения доказывается с помощью принципа
максимума, так же как для первой краевой задачи с разрывными граничными
значениями~[2].

Из единственности следует осевая симметрия~\eqref{eq:axisym}. Коэффициенты
оператора~\eqref{eq:lap-cyl} и граничные условия не зависят от~$\varphi$,
поэтому вместе с функцией $T(r,\varphi,z)$ решением
задачи~\eqref{eq:problem3d} при любом~$\alpha$ является и функция
$T(r,\varphi+\alpha,z)$. Решение единственно, значит, обе функции
совпадают, и $T$ не зависит от~$\varphi$.

\emph{Границы значений температуры.} По принципу максимума~[2,\,3] функция,
удовлетворяющая уравнению Лапласа и отличная от постоянной, не принимает
наибольшего и наименьшего значений во внутренних точках области. На торцах
цилиндра экстремум также невозможен. Условия~\eqref{eq:bc-ends} позволяют
продолжить решение чётным образом через плоскость $z=0$:
$T(r,-z)=T(r,z)$. Продолженная функция удовлетворяет уравнению Лапласа
в цилиндре $-h<z<h$ и принимает те же значения, что и исходная, а точки
торца $z=0$ для неё внутренние. Так же рассматривается торец $z=h$. Поэтому
наибольшее и наименьшее значения температуры определяются её значениями на
боковой поверхности, и при $T_0>0$
\begin{equation}
\label{eq:bounds}
0\le T(r,z)\le T_0 .
\end{equation}
Разрыв функции~$\mu(z)$ оценку~\eqref{eq:bounds} не нарушает: для
ограниченных решений она сохраняется~[2].

Из принципа максимума следует также устойчивость решения: если граничную
функцию~$\mu(z)$ изменить не более чем на~$\varepsilon$, то решение
изменится не более чем на~$\varepsilon$. Существование решения
показано в разделе~\ref{ch:analytic}: решение построено в явном виде
методом функции Грина.

\section{Безразмерная форма задачи}
\label{sec:dimless}

Введём безразмерные переменные
\begin{equation}
\label{eq:dimless-vars}
\xi=\frac{r}{a},\qquad \zeta=\frac{z}{h},\qquad \theta=\frac{T}{T_0} .
\end{equation}
Тогда $\partial/\partial r=(1/a)\,\partial/\partial\xi$,
$\partial/\partial z=(1/h)\,\partial/\partial\zeta$, и после умножения
уравнения~\eqref{eq:main-div} на $a^2/T_0$ задача~\eqref{eq:problem}
принимает вид
\begin{equation}
\label{eq:dimless}
\begin{cases}
\dfrac{1}{\xi}\,\pd{}{\xi}\Bigl(\xi\,\pd{\theta}{\xi}\Bigr)
+\Bigl(\dfrac{a}{h}\Bigr)^{2}\pdd{\theta}{\zeta}=0,
  & 0<\xi<1,\ \ 0<\zeta<1,\\[10pt]
\pd{\theta}{\zeta}\bigg|_{\zeta=0}=0,\qquad
\pd{\theta}{\zeta}\bigg|_{\zeta=1}=0,
  & 0\le\xi<1,\\[10pt]
\theta\big|_{\xi=1}=\mu_0(\zeta),
  & 0<\zeta<1,\\[6pt]
\bigl|\theta\big|_{\xi=0}\bigr|<\infty,
  & 0\le\zeta\le 1,
\end{cases}
\postdisplaypenalty=10000
\end{equation}
где $\mu_0(\zeta)=\mu(h\zeta)/T_0$~— безразмерная температура боковой
поверхности: $\mu_0=0$ при \mbox{$0<\zeta<c/h$} и $\mu_0=1$ при
\mbox{$c/h\le\zeta<1$}.

В задачу~\eqref{eq:dimless} входят только два безразмерных параметра~—
отношение радиуса цилиндра к его высоте $a/h$ и относительная высота линии
скачка температуры $c/h$; они и определяют форму температурного поля.
Температура~$T_0$ служит масштабом:
\begin{equation}
\label{eq:scaling}
T(r,z)=T_0\,\theta\Bigl(\frac{r}{a},\frac{z}{h}\Bigr).
\end{equation}
